\section{The Mexican penitenciary system}
\label{sec:back}

\subsection*{Institutional framework}

In Mexico, criminal jurisdiction is divided between state authorities and the federal government, depending on the nature of the offense. \textbf{State crimes} (\textit{delitos del fuero común}) cover the majority of everyday criminal activity—such as theft, robbery, assault, homicide, and most property or violent crimes—committed within the boundaries of a single state. These cases are investigated by state prosecutors, judged in state courts, and, if custodial sentences are imposed, offenders are generally sent to state prisons. By contrast, \textbf{federal crimes} (\textit{delitos del fuero federal}) are those that either involve federal interests or cross state or national boundaries. They include drug trafficking, organized crime, firearms offenses, money laundering, customs and tax violations, as well as crimes against federal institutions (e.g., corruption of federal officials, attacks on federal infrastructure).  

The legal framework, however, introduces important flexibility in the administration of sentences. \textbf{Article 3 of the Ley que Establece las Normas Mínimas sobre Readaptación Social de Sentenciados} provides that: \textit{“It may also be agreed that inmates sentenced for federal crimes serve their sentences in penitentiary centers under the responsibility of State Governments, when these centers are located closer to their place of residence than those of the Federal Executive.”} This provision highlights the principle that social reintegration is more likely when prisoners remain close to their families and communities, even if this requires federal inmates to serve their sentences in state facilities.  

Similarly, the \textbf{Mexican Constitution (Article 19)} allows for agreements between federal and state authorities so that individuals sentenced for crimes within one jurisdiction may serve their sentences in prisons administered by another. It specifies: \textit{“The Federation, the States, and the Federal District may enter into agreements so that individuals sentenced for crimes under their jurisdiction serve their penalties in penitentiary establishments dependent on a different jurisdiction... Sentenced individuals, under the cases and conditions established by law, may serve their sentences in the penitentiary centers closest to their residence, in order to foster reintegration into the community as a form of social reinsertion.”}\footnote{For a better description see this article by the \href{http://appweb.cndh.org.mx/biblioteca/archivos/pdfs/Clasificacion-penitenciaria.pdf}{CNDH} \citep{villanueva2016clasificacion}.}  

The \textbf{Ley Nacional de Ejecución Penal (Article 49)} further reinforces this principle, stating that: \textit{“Sentenced individuals may serve their sentence in the prison facilities closest to their place of residence.”} An exception is made for those convicted of organized crime, who are confined to specially designated maximum-security facilities such as CEFERESO No. 1 Altiplano in the State of Mexico, opened in 1991.\footnote{This center, as well as all organized crime–related cases, are excluded from the analysis, since they do not meet the institutional feature of geographic proximity necessary for the empirical strategy.}  

In principle, individuals convicted of federal crimes should serve their sentences in federal prisons, while those convicted of state crimes should do so in state facilities. In practice, however, the legal flexibility provided by these rules has produced a recurrent presence of federal prisoners in state prisons. Descriptive evidence indicates that this phenomenon contributes to the chronic overcrowding of local facilities. Data from 2012 shows that approximately 66\% of state and municipal prisons held at least one federal prisoner during the year; in 6\% of these prisons, federal inmates made up more than 30\% of the total population. The most extreme case was the CERESO Calera de Víctor Rosales in Zacatecas, where 72\% of inmates were under federal jurisdiction. 

\subsection*{Pre-trial detention and the reform of 2008}

Understanding the organizational features of the Mexican legal and criminal system is essential to understand the results later provided. In the mixed criminal procedure that prevailed in Mexico prior to the 2008 reform, the process was divided into several stages. It began with a preliminary investigation (\textit{averiguación previa}), conducted by the Public Prosecutor with the assistance of the police and forensic services, aimed at gathering the necessary elements to decide whether to bring charges or dismiss the case. Once the case was presented before a judge, the pre-trial stage (\textit{preinstrucción}) took place, during which the preparatory statement of the accused was heard, and initial evidence could be reviewed. At this point, the judge decided whether to issue an order of processing (\textit{auto de formal prisión} or \textit{auto de sujeción a proceso}) or to release the defendant for lack of evidence. If the case moved forward, it entered the instruction stage (\textit{instrucción}), a phase dedicated to the offering, admission, and presentation of evidence in order to establish the historical truth of the facts. Afterward, the case proceeded to the trial stage (\textit{juicio}), where the parties presented their conclusions, and the judge delivered a ruling in the first instance. Finally, the decision could be challenged in the appeal stage (\textit{segunda instancia}), where higher courts reviewed the application of the law and the evaluation of evidence in the initial proceedings \citep{cossio2015formalprision}. 

Under Mexico’s former mixed criminal procedure system (inquisitorial in nature), criminal proceedings were divided into several stages, with the judge playing a central and active role in the investigation. After the public prosecutor (\textit{Ministerio Público}) presented an initial accusation, the first procedural stage, the pre-trial one, began. At this point, the judge conducted a preliminary evaluation of the evidence presented. Within a maximum of 72 hours from the defendant’s initial appearance, the judge had to decide whether there were sufficient elements to presume the accused’s probable responsibility. If so, the judge issued an \textit{auto de formal prisión} (formal imprisonment order), which both subjected the person to trial and justified their pre-trial detention. This decision had profound implications: once issued, the accused remained detained throughout the instruction and trial stages unless extraordinary remedies succeeded. In other words, pre-trial detention was the default once the \textit{auto de formal prisión} was granted \citep{cossio2015formalprision}. This legal and institutional context allows to test other aspects of judicial decisions that are affected directly by capacity or salience issues, as individuals sent to pre-trial detention occupy a space (separated from already sentenced individuals) on these prisons.

By contrast, following the 2008 constitutional reform and the gradual transition to a new adversarial criminal system (completed nationwide in 2016), this figure was replaced by the \textit{auto de vinculación a proceso} (order linking the accused to criminal proceedings). Under the new system, the judge’s role is limited to guaranteeing procedural rights and overseeing legality, while the prosecutor retains control over the investigation. The linkage to criminal proceedings is also issued within 72 hours but only determines that there is sufficient evidence to initiate a formal trial process — it does not automatically entail pre-trial detention. Instead, preventive detention (\textit{prisión preventiva}) must be requested and justified separately by the prosecution, and the judge must assess whether it is strictly necessary based on specific legal grounds (such as risk of flight or danger to the victim or society) \citep{cossio2015formalprision}.

Although the reform was enacted in 2008, its implementation was gradual and only completed nationwide in 2016. There is no reason to be concerned that this reform biases the results for several reasons. First, the data analyzed in this study only includes cases processed under the mixed criminal procedure system. In fact, up until 2012 the variable that records the procedural status of individuals continued to use the former categories: \textit{auto de formal prisión}, \textit{auto de sujeción a proceso}, release for lack of evidence, and so forth. One potential concern is that the pilot implementation of the new adversarial system might have generated spillover effects on cases still being processed under the old system. However, this possibility is remote, particularly because such pilot programs would have needed to coincide both temporally and geographically with the construction of new federal prisons in order to bias the DID estimates.

\subsection*{Overcapacity over the years}

The Mexican penitentiary system has long suffered significant overextension at both the federal and local levels. Figure \ref{fig:capacity_trend} plots the trend in prison population levels and installed capacity from 2000 to 2012 (period of analysis). The first point to emphasize is that the system has consistently operated overcapacity, with this imbalance widening over time. While the prison population rose steadily, the second important observation is that installed capacity—especially at the local level—also increased, though at a slower pace. This suggests two important dynamics at play during this period: a sustained rise in policing, incarceration rates and/or sentencing length, alongside ongoing though insufficient efforts to expand prison infrastructure at the state level.

In addition to this, during the 2000s, the rise to power of the Partido Acción Nacional (PAN) marked a significant shift in violence and security policy. Often described as a major crime crackdown \citep{DellMexicanDrug, AnnualReviewPoliSci}, PAN governments expanded both federal and local policing and prosecution. Figure \ref{fig:processed} illustrates this pattern. Under Vicente Fox (2001–2006), there was a marked increase in the prosecution of local (state) crimes. In stark contrast, upon taking office, Calderón launched an intense and controversial campaign against drug cartels—widely known as \textit{the war on drugs}. His administration, driven by the aim of curbing the growing influence of organized crime, dramatically expanded law enforcement operations, which led to a sharp rise in incarceration rates at the federal level, prosecuting drug and gun related crimes more actively. By 2010, in Mexico City alone, the number of prisoners had more than doubled, and the local incarceration rate had reached 426 per 100,000 individuals \citep{SSP}. Interestingly, prosecutions of state-level crimes dropped during this period.

% The aggressive prosecution of individuals involved in drug-related activities—often regardless of their direct connection to violent crime—generated a complex and multifaceted socio-political context. The policies enacted during this period had profound implications not only for the criminal justice system but also for the broader social fabric, sparking widespread debate about the effectiveness and consequences of this hardline approach to organized crime \citep{Flores, Atuesta1}.

One immediate consequence of this policy was the continued strain on the federal prison system. As shown in Figure~\ref{fig:processed}, the number of individuals processed for federal offenses increased substantially between 2000 and 2012. In response—partly as a reaction to, and partly in anticipation of, this growth—President Fox initiated the construction of five publicly operated federal prison facilities across different regions of Mexico. This expansion was later continued under President Calderón, who oversaw the construction of three additional public centers.

Subsequently, President Calderón awarded contracts for the construction of ten new Centros Federales de Readaptación Social (CEFERESOS), marking a shift toward private participation in the federal prison system. The policy was justified on the grounds that the private sector possessed greater financial capacity to expand infrastructure, improve prison conditions, and facilitate inmate reintegration. The contracts, granted to a consortium of private firms, covered the construction, maintenance, and operation of each facility \citep{Evalua}. Under this model, private operators were responsible for food provision, facility upkeep, recreational and reintegration programs, and substance abuse treatment, while public authorities retained control over internal security.

To the best of my knowledge, there is limited public documentation detailing President Fox’s rationale for initiating the initial prison expansion. By contrast, President Calderón explicitly articulated the motivation behind the subsequent move toward privatization. The first formal step occurred in 2009, when a presidential order mandated the relocation of all individuals charged with federal offenses to prisons under federal jurisdiction, as many were still being held in state facilities. At the time, however, federal infrastructure was insufficient to implement this mandate. Fully complying with the order would have required the construction of at least eight new prison complexes, yet the government argued that the necessary public resources were unavailable \citep{ZaldivarGonzalez2020Privatizacion}.

This rationale suggests that the policy was driven primarily by federal jurisdictional and institutional constraints rather than by localized trends in state-level prison overcrowding. The central issue was not uneven congestion across regions, but rather the lack of adequate federal infrastructure to comply with existing legal mandates governing the incarceration of individuals convicted of federal crimes.



